problem:  
Let \((\mathbb{R}^2,d_{SR},\mu_\Phi)\) be the **weighted Grushin plane**, where the horizontal distribution is spanned by the orthonormal frame
\[
X_1 = \partial_x, \qquad X_2 = |x|^{\alpha}\partial_y, \quad \alpha > 0,
\]
the Carnot–Carathéodory distance \(d_{SR}\) is induced by this frame, and  
\(\mu_\Phi = e^{-\Phi(x,y)}\,dx\,dy\) is a smooth weighted measure with potential \(\Phi \in C^\infty_c(\mathbb{R}^2)\).  

**Conjecture (Weighted rigidity of synthetic curvature on the Grushin plane):**  
For every \(\alpha>0\) and every smooth potential \(\Phi\), the metric‑measure space \((\mathbb{R}^2,d_{SR},\mu_\Phi)\) **fails** to satisfy the Lott–Sturm–Villani curvature–dimension condition \(\mathrm{CD}(K,N)\) for any \(K\in\mathbb{R}, N>1\).

Equivalently, synthetic Ricci lower bounds in the LSV sense are possible **only** when \(\alpha=0\) (the Euclidean plane case).

A complementary possibility—if the conjecture fails—is to characterize explicitly all pairs \((\alpha,\Phi)\) for which \(\mathrm{CD}(K,N)\) might hold, thereby generating the first known examples of genuinely degenerate sub‑Riemannian geometries supporting synthetic curvature bounds.

Why is it a "good" problem:  
This formulation is now maximally concrete and logically consistent. It focuses on a single, canonical non‑equiregular model where horizontal directions vanish along a line, leading to variable step and branching geodesics. By asking whether **any** weighting can salvage a Lott–Sturm–Villani curvature–dimension inequality, it strikes directly at the core of the compatibility problem between synthetic curvature and degenerate sub‑Riemannian geometry.  

A resolution—positive or negative—would clarify whether metric‑measure curvature bounds can ever coexist with horizontal degeneracy, determining the true boundaries of the \(\mathrm{CD}(K,N)\) theory. The question is simple to state using explicit vector fields and familiar weights but conceptually deep, engaging analytic, geometric, and optimal‑transport methods. It thereby exemplifies a “narrow‑entrance, profound‑depth” research problem bridging Riemannian and sub‑Riemannian curvature theories.